\documentclass[letterpaper,12pt,oneside]{article}

% --- CONFIGURACIÓN DE PÁGINA (GEOMETRÍA) ---
\usepackage[
    left=25mm, 
    right=25mm,
    top=30mm,
    bottom=30mm,
    headheight=15pt, 
]{geometry}

% --- PAQUETES DE IDIOMA Y FUENTES ---
\usepackage[T1]{fontenc}
\usepackage[utf8]{inputenc}
\usepackage[spanish,es-nodecimaldot,es-tabla]{babel}

% --- PAQUETES GRÁFICOS Y DE DISEÑO ---
\usepackage{xcolor}
\usepackage{tcolorbox}
\tcbuselibrary{listings, fitting}

% --- HIPERVÍNCULOS ---
\usepackage{hyperref}
\hypersetup{
    colorlinks=true,
    linkcolor=blue,
    citecolor=blue,
    urlcolor=blue
}

% --- COLORES PARA CAJAS DE CÓDIGO ---
\definecolor{bgcode}{HTML}{F4F4F4}
\definecolor{framecode}{HTML}{4A4A4A}

% --- ENTORNO PARA COMANDOS GIT ---
\newtcblisting{gitbox}[1][]{
    colback=bgcode,
    colframe=framecode,
    arc=3pt,
    boxrule=1pt,
    left=5pt,
    right=5pt,
    top=5pt,
    bottom=5pt,
    fontupper=\ttfamily\small,
    listing only,
    title=#1,
    coltitle=white,
    fonttitle=\bfseries
}

% --- ENCABEZADOS ---
\usepackage{fancyhdr}
\pagestyle{fancy}
\fancyhf{}
\fancyhead[C]{\textbf{Guía Completa de Git y Trabajo Colaborativo}}
\fancyfoot[C]{\thepage}
\renewcommand{\headrulewidth}{0.4pt}
\renewcommand{\footrulewidth}{0.4pt}

\begin{document}

\begin{center}
    {\Huge \textbf{Guía de Configuración y Uso de Git}}\\
    \vspace{0.5cm}
    {\large \textit{Manual integral para trabajo colaborativo en repositorio}}
\end{center}

\vspace{1cm}

\section{Configuración Inicial del Colaborador}
Antes de comenzar a trabajar en cualquier proyecto usando Git, cada colaborador debe configurar su identidad. Esto es crucial, ya que los *commits* (guardados) llevarán este nombre y correo, permitiendo saber quién hizo qué cambio.

\begin{gitbox}[Configuración de Identidad]
git config --global user.name "Tu Nombre Completo"
git config --global user.email "tu_correo@ejemplo.com"
\end{gitbox}

Para verificar que los datos se han guardado correctamente, puedes listar tu configuración con:
\begin{gitbox}[]
git config --list
\end{gitbox}

\section{Obtener el Repositorio}
Para comenzar a trabajar, necesitas descargar el proyecto desde la nube (GitHub, GitLab, etc.) a tu computadora local. Esto se hace una sola vez.

\begin{gitbox}[Clonar el repositorio]
git clone https://ruta-del-repositorio.git
\end{gitbox}

\section{Manejo de Ramas (Branches)}
Las ramas permiten a los colaboradores trabajar en nuevas características sin afectar el código principal (usualmente llamado \texttt{main} o \texttt{master}).

\subsection{Crear y listar ramas}
Para ver en qué rama te encuentras actualmente y la lista de ramas locales:
\begin{gitbox}[]
git branch
\end{gitbox}

Para crear una rama nueva y moverte a ella inmediatamente:
\begin{gitbox}[Crear rama y cambiar a ella]
git checkout -b nombre_de_tu_rama
# Alternativa en versiones modernas de Git:
git switch -c nombre_de_tu_rama
\end{gitbox}

\subsection{Moverse entre ramas}
Si necesitas saltar de la rama en la que estás hacia otra (por ejemplo, para revisar cómo iba \texttt{main}), asegúrate primero de no tener cambios pendientes.
\begin{gitbox}[Navegación entre ramas]
git checkout nombre_de_la_rama_existente
# Alternativa moderna:
git switch nombre_de_la_rama_existente
\end{gitbox}

\section{Flujo de Trabajo (Guardar y Subir Cambios)}
Una vez que has modificado archivos, el flujo tradicional consiste en 3 pasos: preparar (\textit{add}), empaquetar (\textit{commit}) y publicar (\textit{push}).

\textbf{Paso 1: Añadir los cambios.} 
\begin{gitbox}[Preparar archivos]
git add .
# El punto (.) añade todos los archivos modificados. 
# Si solo quieres añadir uno: git add nombre_archivo.tex
\end{gitbox}

\textbf{Paso 2: Guardar el estado (Commit).}
\begin{gitbox}[Crear el Commit]
git commit -m "Mensaje descriptivo sobre lo que se modificó"
\end{gitbox}

\textbf{Paso 3: Subir a la nube (Push).}
\begin{gitbox}[Subir los cambios]
git push origin nombre_de_tu_rama
# Si es la primera vez que subes esa rama, usa la bandera -u:
git push -u origin nombre_de_tu_rama
\end{gitbox}

\textbf{Descargar cambios del equipo (Pull).}
Siempre antes de empezar a programar en el día, es buena idea traer los últimos avances de tus compañeros.
\begin{gitbox}[Actualizar localmente]
git pull origin main
\end{gitbox}

\section{Resolución de Errores y Comandos de Auxilio}
Es normal equivocarse. Git tiene comandos valiosos para saber qué está pasando o para retroceder en el tiempo.

\subsection{Ver el estado actual}
Si no sabes qué archivos modificaste o dónde estás parado, este es tu mejor amigo.
\begin{gitbox}[Revisar estado]
git status
\end{gitbox}

\subsection{Historial de cambios (\textit{Logs})}
Para ver quién hizo los últimos commits y cuál es su ID.
\begin{gitbox}[]
git log --oneline
\end{gitbox}

\subsection{Deshacer cosas}
\textbf{1. Quitar algo del \texttt{git add} previo al \texttt{commit}:}
Se te fue un archivo basura por accidente:
\begin{gitbox}[Restaurar stage]
git restore --staged archivo_equivocado.tex
\end{gitbox}

\textbf{2. Borrar las últimas modificaciones de un archivo:}
ATENCIÓN: Esto borrará tus cambios y lo dejará como estaba en el último commit.
\begin{gitbox}[Limpiar archivo local]
git restore archivo_que_quieres_regresar.tex
\end{gitbox}

\textbf{3. Equivocación en el último \texttt{commit} (pero sin haber hecho \texttt{push}):}
Si enviaste un commit incompleto o escribiste mal el mensaje:
\begin{gitbox}[Deshacer último commit]
git reset HEAD~1
\end{gitbox}

\end{document}